\documentclass[a4paper,12pt]{article}

% ---------------- PACOTES ----------------
\usepackage[utf8]{inputenc}
\usepackage[T1]{fontenc}
\usepackage[brazil]{babel}
\usepackage{geometry}
\usepackage{graphicx}
\usepackage{amsmath, amssymb}
\usepackage{siunitx}
\usepackage{caption}
\usepackage{subcaption}
\usepackage{float}
\usepackage{indentfirst}
\usepackage{hyperref}
\usepackage{array}

\geometry{margin=2.5cm}

% ---------------- INÍCIO ----------------
\begin{document}

% -------- CAPA --------
\begin{titlepage}

    % Logo
    \noindent
    \begin{minipage}[c]{0.18\textwidth}
        \includegraphics[width=\linewidth]{png-logo-ufma-colorido.png}
        \end{minipage}
    \hfill
    \begin{minipage}[c]{0.78\textwidth}
    \raggedright
    {\small \textbf{UNIVERSIDADE FEDERAL DO MARANHÃO - UFMA\\
    CENTRO DE CIÊNCIA EXATAS E TECNOLOGIA - CCET\\
    ENGENHARIA DA COMPUTAÇÃO - EECP\\
    ARQUITETURA DE COMPUTADORES
    }}
    \end{minipage}

    \vspace{2cm}

    % Título
    \begin{center}
    {\Large \textbf{Sistema Embarcado de Monitoramento Ambiental e Meteorológico}}
    \end{center}

    \vspace{2cm}

    % Discentes
    \begin{center}
    {\large Discentes:\\
    \large \textbf{}\\
    \large \textbf{João Henrique Calixto Teixeira (2023030390)}\\
    \large \textbf{Marcos Paulo Dominice Silva (2023030318)}\\
    \large \textbf{Rai Abreu Machado (2023034498)}\\
    \large \textbf{Ronald Loureiro Camara (2023034460)}\\
    \large \textbf{Wesley Ferreira Costa (2023034442)}}\\
    \end{center}

    \vspace{1.5cm}

    % Docente
    \begin{center}
    {\large Docente:\\
    \textbf{Luiz Henrique Neves Rodrigues}}
    \end{center}

    \vfill

    % Rodapé
    \begin{center}
    {\large \textbf{São Luís - MA}}\\
    {\large \textbf{2026}}
    \end{center}

\end{titlepage}

% -------- SUMÁRIO --------
\tableofcontents
\newpage
% -------- SEÇÕES --------

\section{Introdução}
A arquitetura de Von Neumann constitui uma das bases fundamentais da computação moderna, estabelecendo a organização clássica de sistemas computacionais através da divisão entre entrada de dados, processamento, memória e saída de informações. Mesmo após décadas de sua proposição, seus princípios continuam presentes em computadores modernos, microcontroladores e sistemas embarcados.

Nesse contexto, o presente trabalho apresenta o desenvolvimento de um sistema embarcado de monitoramento ambiental utilizando o Arduino Nano, baseado no microcontrolador ATmega328P, com o objetivo principal de demonstrar na prática a aplicação dos conceitos da arquitetura clássica de Von Neumann em um sistema computacional real.

O projeto foi desenvolvido utilizando sensores responsáveis pela aquisição de dados ambientais, incluindo temperatura, umidade relativa do ar, pressão atmosférica, altitude aproximada, luminosidade e qualidade relativa do ar. Para isso, foram utilizados os sensores BME280, BH1750 e MQ-135, integrados ao microcontrolador através de diferentes formas de comunicação e aquisição de dados.

Os sensores digitais BME280 e BH1750 utilizam o protocolo I2C para transferência de informações ao microcontrolador, enquanto o sensor MQ-135 opera através de um sinal analógico variável, exigindo a utilização do conversor analógico-digital (ADC) interno do ATmega328P. Dessa forma, o projeto também demonstra o funcionamento do processo de conversão de sinais analógicos em dados digitais dentro de um sistema computacional embarcado.

Durante a execução do sistema, os sensores atuam como dispositivos de entrada, fornecendo dados ambientais ao microcontrolador. O ATmega328P realiza o processamento dessas informações através do firmware desenvolvido em linguagem C/C++, utilizando sua unidade de processamento e memória interna para interpretar os dados recebidos. Após o processamento, as informações são enviadas ao Monitor Serial, caracterizando a etapa de saída de dados do sistema.

Além da aquisição e exibição das informações ambientais, o projeto permite observar, na prática, diversos elementos presentes na arquitetura clássica de Von Neumann, como:

\begin{itemize}
    \item Entrada de dados através dos sensores;
    \item Processamento centralizado pelo microcontrolador;
    \item Utilização de memória para armazenamento de instruções e variáveis;
    \item Comunicação entre periféricos e processador;
    \item Saída de informações processadas.
\end{itemize}

Assim, o sistema desenvolvido serve como uma aplicação prática dos conceitos estudados em arquitetura de computadores, permitindo relacionar os fundamentos teóricos da arquitetura de Von Neumann com a implementação real de um sistema embarcado baseado em aquisição, processamento e exibição de dados ambientais.

\section{Objetivos}
O presente projeto tem como objetivo desenvolver um sistema embarcado de monitoramento ambiental utilizando o Arduino Nano baseado no microcontrolador ATmega328P, visando demonstrar, na prática, a aplicação dos conceitos de arquitetura de computadores e realizar uma análise comparativa entre a arquitetura utilizada pelo microcontrolador e a arquitetura clássica de Von Neumann.

Para isso, o sistema foi projetado para realizar a aquisição e processamento de dados ambientais por meio dos sensores BME280, BH1750 e MQ-135, permitindo a coleta de informações relacionadas à temperatura, umidade, pressão atmosférica, luminosidade e qualidade relativa do ar.

Além do desenvolvimento do sistema funcional, o projeto busca analisar aspectos relacionados ao funcionamento interno do microcontrolador, incluindo comunicação com dispositivos de entrada e saída, processamento de dados, utilização de memória e conversão analógico-digital, estabelecendo relações entre os conceitos teóricos estudados em arquitetura de computadores e sua aplicação prática em sistemas embarcados.

\section{Referencial Teórico}
O desenvolvimento do presente projeto fundamenta-se em conceitos relacionados à arquitetura de computadores, sistemas embarcados, aquisição de dados, comunicação serial e processamento digital de sinais. Para a implementação do sistema de monitoramento ambiental, foram utilizados como principais referências técnicas os datasheets do microcontrolador ATmega328P e dos sensores BME280, BH1750 e MQ-135.

\subsection{Arquitetura de Von Neumann}

A arquitetura de Von Neumann constitui um dos principais modelos computacionais utilizados na construção de sistemas digitais modernos. Esse modelo define uma estrutura composta por:
\begin{itemize}
    \item Unidade Central de Processamento (CPU);
    \item Memória;
    \item Dispositivos de entrada;
    \item Dispositivos de saída;
    \item Barramentos de comunicação.
\end{itemize}

Na arquitetura proposta por Von Neumann, tanto dados quanto instruções compartilham o mesmo espaço de memória e são processados sequencialmente pela CPU.

No projeto desenvolvido, os sensores atuam como dispositivos de entrada de dados, enquanto o ATmega328P realiza o processamento das informações adquiridas. A memória interna do microcontrolador armazena as instruções do firmware e as variáveis utilizadas durante a execução do sistema, enquanto o Monitor Serial atua como interface de saída de dados.

Assim, o sistema desenvolvido permite observar, na prática, os principais elementos presentes na arquitetura clássica de Von Neumann aplicados em um sistema embarcado moderno.

\subsection{Microcontrolador ATmega328P}

O ATmega328P é um microcontrolador de 8 bits baseado na arquitetura AVR RISC, desenvolvido pela Atmel. Segundo o datasheet do fabricante, o dispositivo possui arquitetura RISC avançada, 32 registradores de propósito geral e capacidade de execução próxima de 1 MIPS por MHz.

O microcontrolador utilizado no projeto encontra-se integrado à plataforma Arduino Nano, responsável pelo gerenciamento dos sensores e pelo processamento dos dados ambientais coletados.

De acordo com o datasheet, o ATmega328P possui:
\begin{itemize}
    \item 32 KB de memória Flash;
    \item 2 KB de SRAM;
    \item 1 KB de EEPROM;
    \item Conversor analógico-digital (ADC) de 10 bits;
    \item Interface serial I2C (TWI);
    \item Interface USART;
    \item Temporizadores internos;
    \item 23 linhas de entrada e saída digitais.
\end{itemize}

No sistema desenvolvido, o ATmega328P é responsável por:
\begin{itemize}
    \item Realizar aquisição de dados ambientais;
    \item Executar comunicação I2C;
    \item Executar conversão analógico-digital;
    \item Processar os dados dos sensores;
    \item Classificar a qualidade relativa do ar;
    \item Exibir informações no Monitor Serial.
\end{itemize}

\subsection{Comunicação I2C}

O protocolo I2C (\textit{Inter-Integrated Circuit}) é um barramento serial síncrono utilizado para comunicação entre dispositivos eletrônicos.

Esse protocolo utiliza apenas duas linhas de comunicação:
\begin{itemize}
    \item SDA (\textit{Serial Data});
    \item SCL (\textit{Serial Clock}).
\end{itemize}

Segundo o datasheet do ATmega328P, o microcontrolador possui interface TWI compatível com o protocolo I2C.

No projeto desenvolvido, os sensores BME280 e BH1750 compartilham o mesmo barramento I2C do Arduino Nano, utilizando diferentes endereços lógicos para comunicação com o microcontrolador.

A utilização do protocolo I2C reduz a quantidade de conexões físicas necessárias no sistema e demonstra a comunicação entre múltiplos periféricos digitais em sistemas embarcados.

\subsection{Sensor BME280}

O BME280 é um sensor digital desenvolvido pela Bosch Sensortec, capaz de realizar medições de:
\begin{itemize}
    \item Temperatura;
    \item Umidade relativa do ar;
    \item Pressão atmosférica.
\end{itemize}

A partir da pressão atmosférica, também é possível estimar a altitude aproximada.

Segundo o datasheet do fabricante, o BME280 realiza internamente ciclos de medição compostos por:
\begin{enumerate}
    \item Medição de temperatura;
    \item Medição de pressão;
    \item Medição de umidade;
    \item Filtragem opcional dos dados.
\end{enumerate}

O sensor permite utilização de técnicas de \textit{oversampling} para redução de ruído nas medições. Além disso, possui filtro IIR interno para suavização de variações rápidas de pressão e temperatura.

No sistema desenvolvido, o BME280 opera realizando medições contínuas e disponibilizando os dados através da interface I2C.

A utilização do BME280 demonstra a integração entre sensores digitais inteligentes e sistemas embarcados.

\subsection{Sensor BH1750}

O BH1750 é um sensor digital de luminosidade responsável pela medição da intensidade luminosa ambiente em lux.

O sensor utiliza comunicação I2C para envio das informações ao microcontrolador, compartilhando o barramento com o BME280.

Internamente, o BH1750 realiza a conversão da intensidade luminosa em dados digitais, permitindo que o ATmega328P apenas solicite os valores medidos através da interface serial.

No projeto desenvolvido, o BH1750 representa um periférico digital dedicado à aquisição de luminosidade ambiente.

\subsection{Sensor MQ-135}

O MQ-135 é um sensor analógico utilizado para monitoramento relativo da qualidade do ar.

Segundo o datasheet do fabricante, o sensor é capaz de detectar gases como:
\begin{itemize}
    \item NH$_3$;
    \item NOx;
    \item Álcool;
    \item Benzeno;
    \item Fumaça;
    \item CO$_2$.
\end{itemize}

O sensor possui:
\begin{itemize}
    \item Camada semicondutora de dióxido de estanho (SnO$_2$);
    \item Aquecedor interno;
    \item Circuito resistivo sensível a gases.
\end{itemize}

Seu funcionamento baseia-se na variação da resistência elétrica interna do material semicondutor quando exposto a diferentes concentrações de gases.

De acordo com o datasheet:
\begin{itemize}
    \item a resistência do aquecedor é aproximadamente 33$\Omega$;
    \item o sensor opera com alimentação de 5V;
    \item o tempo recomendado de pré-aquecimento é superior a 24 horas.
\end{itemize}

O MQ-135 fornece um sinal analógico contínuo ao microcontrolador, sendo necessária a utilização do ADC interno do ATmega328P para interpretação dos dados.

O datasheet também apresenta curvas de sensibilidade baseadas na razão:

\[
\frac{R_s}{R_o}
\]

Onde:
\begin{itemize}
    \item $R_s$ representa a resistência atual do sensor;
    \item $R_o$ representa a resistência do sensor em ambiente calibrado.
\end{itemize}

Essas curvas permitem estimar concentrações relativas de gases em ppm.

\subsection{Conversão Analógico-Digital (ADC)}

O conversor analógico-digital (ADC) é responsável por transformar sinais analógicos contínuos em valores digitais processáveis pelo sistema computacional.

Segundo o datasheet do ATmega328P, o microcontrolador possui ADC de 10 bits com até 8 canais analógicos.

No Arduino Nano, o MQ-135 encontra-se conectado ao pino A0, fornecendo uma tensão analógica variável proporcional à concentração relativa de gases.

O ADC converte tensões entre 0V e 5V em valores digitais entre 0 e 1023.

A conversão pode ser representada por:

\[
ADC = \frac{V_{in}}{V_{ref}} \times 1023
\]

Onde:
\begin{itemize}
    \item $V_{in}$ representa a tensão analógica aplicada na entrada;
    \item $V_{ref}$ representa a tensão de referência do ADC.
\end{itemize}

Como o ADC possui resolução de 10 bits:

\[
2^{10} = 1024
\]

Assim:
\begin{itemize}
    \item 0V corresponde a 0;
    \item 5V corresponde a 1023.
\end{itemize}

No firmware desenvolvido, a leitura é realizada através da função:

\begin{verbatim}
analogRead(MQ135_PIN);
\end{verbatim}

O valor digital obtido é utilizado pelo sistema para classificar a qualidade relativa do ar em diferentes níveis.

\subsection{Sistemas Embarcados}

Sistemas embarcados são sistemas computacionais desenvolvidos para executar funções específicas dentro de dispositivos eletrônicos.

Diferentemente de computadores de propósito geral, sistemas embarcados possuem:
\begin{itemize}
    \item Hardware dedicado;
    \item Baixo consumo energético;
    \item Processamento específico;
    \item Integração entre software e hardware.
\end{itemize}

No projeto desenvolvido, o Arduino Nano atua como sistema embarcado responsável pelo gerenciamento dos sensores, aquisição dos dados ambientais e processamento das informações em tempo real.

O sistema implementado demonstra conceitos fundamentais de:
\begin{itemize}
    \item integração hardware/software;
    \item aquisição de dados;
    \item processamento digital;
    \item comunicação entre periféricos;
    \item arquitetura de computadores.
\end{itemize}

Além disso, o projeto evidencia a aplicação prática da arquitetura de Von Neumann em sistemas embarcados modernos.

\section{Materiais e Metodologia}

\subsection{Materiais}

\begin{enumerate}

    \item \textbf{Arduino Nano} -- Plataforma de desenvolvimento baseada no microcontrolador ATmega328P, responsável pelo processamento dos dados, gerenciamento dos sensores e execução do firmware do sistema embarcado.
    
    \item \textbf{Protoboard} -- Placa de prototipagem utilizada para montagem do circuito eletrônico sem a necessidade de soldagem, permitindo conexões rápidas entre os componentes do sistema.
    
    \item \textbf{Sensor BME280} -- Sensor digital utilizado para medição de temperatura, umidade relativa do ar e pressão atmosférica, além da estimativa aproximada de altitude.
    
    \item \textbf{Sensor BH1750} -- Sensor digital utilizado para medição da intensidade luminosa ambiente em lux.
    
    \item \textbf{Sensor MQ-135} -- Sensor analógico utilizado para monitoramento relativo da qualidade do ar através da detecção de gases e compostos voláteis presentes no ambiente.
    
    \item \textbf{Jumpers} -- Fios condutores utilizados para realizar as conexões elétricas entre o Arduino Nano, sensores e protoboard.
    
    \item \textbf{Cabo USB} -- Utilizado para alimentação do Arduino Nano e comunicação serial com o computador durante a programação e monitoramento dos dados.
    
    \item \textbf{Computador} -- Equipamento utilizado para desenvolvimento do firmware, gravação do código no microcontrolador e visualização dos dados através do Monitor Serial da IDE Arduino.
    
    \item \textbf{IDE Arduino} -- Ambiente de desenvolvimento utilizado para escrita, compilação e upload do firmware para o microcontrolador ATmega328P.
    
\end{enumerate}
\subsection{Metodologia}

Inicialmente, foi realizada a montagem do circuito eletrônico em protoboard, integrando os sensores BME280, BH1750 e MQ-135 ao sistema embarcado.

Os sensores BME280 e BH1750 foram conectados utilizando o protocolo de comunicação I2C, compartilhando as linhas SDA e SCL do microcontrolador. Já o sensor MQ-135 foi integrado através de sua saída analógica, conectada ao pino A0 para utilização do conversor analógico-digital (ADC).

A Tabela \ref{tab:bme280} apresenta as conexões utilizadas no sensor BME280.

\begin{table}[H]
\centering
\caption{Conexões do sensor BME280}
\label{tab:bme280}
\begin{tabular}{|c|c|}
\hline
\textbf{BME280} & \textbf{Conexão} \\ \hline
VIN & 3V3 \\ \hline
GND & GND \\ \hline
SDA & A4 \\ \hline
SCL & A5 \\ \hline
\end{tabular}
\end{table}

A Tabela \ref{tab:bh1750} apresenta as conexões utilizadas no sensor BH1750.

\begin{table}[H]
\centering
\caption{Conexões do sensor BH1750}
\label{tab:bh1750}
\begin{tabular}{|c|c|}
\hline
\textbf{BH1750} & \textbf{Conexão} \\ \hline
VCC & 3V3 \\ \hline
GND & GND \\ \hline
SDA & A4 \\ \hline
SCL & A5 \\ \hline
ADDR & GND \\ \hline
\end{tabular}
\end{table}

O sensor MQ-135 foi conectado conforme apresentado na Tabela \ref{tab:mq135}.

\begin{table}[H]
\centering
\caption{Conexões do sensor MQ-135}
\label{tab:mq135}
\begin{tabular}{|c|c|}
\hline
\textbf{MQ-135} & \textbf{Conexão} \\ \hline
VCC & 5V \\ \hline
GND & GND \\ \hline
AO & A0 \\ \hline
\end{tabular}
\end{table}

Após a montagem física do circuito, foi desenvolvido o firmware responsável pela inicialização dos sensores, aquisição dos dados ambientais, processamento das informações e exibição dos resultados no Monitor Serial.

Para comunicação com os sensores digitais, foram utilizadas as bibliotecas:
\begin{itemize}
    \item \texttt{Wire.h};
    \item \texttt{Adafruit\_Sensor.h};
    \item \texttt{Adafruit\_BME280.h};
    \item \texttt{BH1750.h}.
\end{itemize}

O sensor BME280 foi inicializado utilizando o endereço I2C \texttt{0x76}, enquanto o BH1750 operou utilizando seu endereço padrão \texttt{0x23}.

Durante a execução do firmware, o sistema realiza leituras periódicas de:
\begin{itemize}
    \item temperatura;
    \item umidade relativa do ar;
    \item pressão atmosférica;
    \item altitude aproximada;
    \item luminosidade ambiente;
    \item qualidade relativa do ar.
\end{itemize}

A leitura do sensor MQ-135 foi realizada através da função:

\begin{verbatim}
analogRead(MQ135_PIN);
\end{verbatim}

O valor analógico fornecido pelo sensor é convertido pelo ADC de 10 bits do microcontrolador em um valor digital entre 0 e 1023.

Após a conversão analógico-digital, o firmware realiza uma classificação relativa da qualidade do ar conforme intervalos previamente definidos no código.

A Tabela \ref{tab:qualidade_ar} apresenta os níveis utilizados no sistema.

\begin{table}[H]
\centering
\caption{Classificação relativa da qualidade do ar}
\label{tab:qualidade_ar}
\begin{tabular}{|c|c|}
\hline
\textbf{Faixa ADC} & \textbf{Classificação} \\ \hline
0 -- 250 & Boa \\ \hline
251 -- 450 & Moderada \\ \hline
451 -- 700 & Ruim \\ \hline
Acima de 700 & Crítica \\ \hline
\end{tabular}
\end{table}

O firmware final utilizado no projeto é apresentado a seguir.

\begin{verbatim}
#include <Wire.h>
#include <Adafruit_Sensor.h>
#include <Adafruit_BME280.h>
#include <BH1750.h>

// =========================
// OBJETOS DOS SENSORES
// =========================

Adafruit_BME280 bme;
BH1750 lightMeter;

// =========================
// CONFIGURAÇÕES
// =========================

#define SEALEVELPRESSURE_HPA (1013.25)

// Pino analógico do MQ-135
#define MQ135_PIN A0

// Variável para verificar se o BME280 foi encontrado
bool bmeStatus = false;


// =========================
// FUNÇÃO: CLASSIFICAÇÃO DO AR
// =========================

String classificarQualidadeAr(int valor) {

  if (valor <= 250) {

    return "BOA";

  }

  else if (valor <= 450) {

    return "MODERADA";

  }

  else if (valor <= 700) {

    return "RUIM";

  }

  else {

    return "CRITICA";

  }

}


// =========================
// SETUP
// =========================

void setup() {

  Serial.begin(9600);
  Wire.begin();

  Serial.println("====================================");
  Serial.println(" SISTEMA DE MONITORAMENTO AMBIENTAL ");
  Serial.println("====================================");

  // =========================
  // Inicialização BME280
  // =========================

  bmeStatus = bme.begin(0x76);

  if (!bmeStatus) {

    Serial.println("Erro ao encontrar o BME280!");
    Serial.println("Verifique conexoes e endereco I2C.");
    Serial.println();

  }

  else {

    Serial.println("BME280 inicializado com sucesso!");

  }

  // =========================
  // Inicialização BH1750
  // =========================

  lightMeter.begin();

  Serial.println("BH1750 inicializado com sucesso!");

  // =========================
  // Inicialização MQ-135
  // =========================

  pinMode(MQ135_PIN, INPUT);

  Serial.println("MQ-135 inicializado com sucesso!");

  Serial.println();

}


// =========================
// LOOP PRINCIPAL
// =========================

void loop() {

  Serial.println("========== DADOS AMBIENTAIS ==========");
  Serial.println();

  // =========================
  // BME280
  // =========================

  if (bmeStatus) {

    float temperatura = bme.readTemperature();
    float umidade = bme.readHumidity();
    float pressao = bme.readPressure() / 100.0F;
    float altitude = bme.readAltitude(SEALEVELPRESSURE_HPA);

    Serial.print("Temperatura: ");
    Serial.print(temperatura);
    Serial.println(" C");

    Serial.print("Umidade: ");
    Serial.print(umidade);
    Serial.println(" %");

    Serial.print("Pressao: ");
    Serial.print(pressao);
    Serial.println(" hPa");

    Serial.print("Altitude Aproximada: ");
    Serial.print(altitude);
    Serial.println(" m");

  }

  else {

    Serial.println("BME280 indisponivel.");

  }

  Serial.println();

  // =========================
  // BH1750
  // =========================

  float luminosidade = lightMeter.readLightLevel();

  Serial.print("Luminosidade: ");
  Serial.print(luminosidade);
  Serial.println(" lux");

  Serial.println();

  // =========================
  // MQ-135
  // =========================

  int qualidadeAr = analogRead(MQ135_PIN);

  Serial.print("Qualidade do Ar: ");
  Serial.print(qualidadeAr);

  Serial.print("  [");

  Serial.print(classificarQualidadeAr(qualidadeAr));

  Serial.println("]");

  Serial.println();

  Serial.println("======================================");
  Serial.println();

  delay(2000);

}
\end{verbatim}

Após o processamento das informações ambientais, os dados são exibidos em tempo real através do Monitor Serial da IDE Arduino.

A saída serial utilizada pelo sistema é apresentada a seguir.

\begin{verbatim}
========== DADOS AMBIENTAIS ==========

Temperatura: 27.4 C
Umidade: 63.2 %
Pressao: 1008.4 hPa
Altitude Aproximada: 42.7 m

Luminosidade: 315.0 lux

Qualidade do Ar: 412 [MODERADA]

======================================
\end{verbatim}

Após a implementação do firmware, foram realizados testes individuais e integrados dos sensores, verificando o correto funcionamento da comunicação I2C, da leitura analógica do MQ-135 e da exibição das informações ambientais no Monitor Serial.

\section{Resultados e Discussão}

Após a implementação completa do sistema embarcado, foi realizado um processo contínuo de aquisição de dados ambientais durante um intervalo de 1 hora, compreendido entre 16:30h e 17:30h.

Durante esse período, o firmware executou leituras periódicas dos sensores BME280, BH1750 e MQ-135, registrando informações referentes à:
\begin{itemize}
    \item temperatura;
    \item umidade relativa do ar;
    \item pressão atmosférica;
    \item luminosidade ambiente;
    \item qualidade relativa do ar.
\end{itemize}

Os dados coletados foram armazenados através do Monitor Serial da IDE Arduino e posteriormente exportados para um arquivo de texto (\texttt{.txt}). Após a aquisição, os dados foram processados utilizando a linguagem Python e a biblioteca \texttt{Matplotlib}, responsável pela geração dos gráficos temporais utilizados na análise experimental.

Para construção dos gráficos, foi utilizado um vetor temporal contendo 15 amostras distribuídas uniformemente ao longo de 60 minutos. O módulo \texttt{datetime} foi utilizado para associação dos horários às leituras registradas, permitindo representar a evolução temporal das variáveis ambientais monitoradas.

\subsection{Análise da Temperatura}

A Figura \ref{fig:temperatura} apresenta o comportamento da temperatura ambiente durante o período de monitoramento.

\begin{figure}[H]
    \centering
    \includegraphics[width=0.85\linewidth]{temperatura.png}
    \caption{Variação da temperatura entre 16:30h e 17:30h}
    \label{fig:temperatura}
\end{figure}

Observa-se um crescimento gradual da temperatura ao longo do tempo, partindo de aproximadamente 28,4°C e atingindo valores próximos de 29,3°C.

Esse comportamento pode estar associado ao aquecimento natural do ambiente e também à dissipação térmica produzida pelos próprios componentes eletrônicos durante o funcionamento contínuo do sistema.

As pequenas oscilações observadas demonstram a estabilidade do sensor BME280 e a consistência das leituras realizadas através da comunicação I2C.

\subsection{Análise da Umidade Relativa do Ar}

A Figura \ref{fig:umidade} apresenta a variação da umidade relativa do ar durante o experimento.

\begin{figure}[H]
    \centering
    \includegraphics[width=0.85\linewidth]{umidade.png}
    \caption{Variação da umidade relativa do ar entre 16:30h e 17:30h}
    \label{fig:umidade}
\end{figure}

Os resultados demonstram uma redução gradual da umidade relativa ao longo do intervalo monitorado, variando de aproximadamente 76\% para valores próximos de 72\%.

Esse comportamento apresenta relação direta com o aumento da temperatura observado anteriormente, uma vez que temperaturas mais elevadas reduzem a umidade relativa do ar.

A estabilidade da curva obtida demonstra o correto funcionamento do sistema de aquisição de dados implementado no firmware.

\subsection{Análise da Pressão Atmosférica}

A Figura \ref{fig:pressao} apresenta os valores de pressão atmosférica registrados pelo sistema.

\begin{figure}[H]
    \centering
    \includegraphics[width=0.85\linewidth]{pressão_atmosférica.png}
    \caption{Variação da pressão atmosférica entre 16:30h e 17:30h}
    \label{fig:pressao}
\end{figure}

Os dados indicam pequenas variações em torno de 1008 hPa e 1009 hPa, comportamento esperado para intervalos curtos de monitoramento.

A baixa variação observada evidencia a estabilidade do sensor BME280 e demonstra a capacidade do sistema em realizar medições contínuas de grandezas ambientais com baixa interferência.

\subsection{Análise da Luminosidade}

A Figura \ref{fig:luminosidade} apresenta a variação da luminosidade ambiente registrada pelo sensor BH1750.

\begin{figure}[H]
    \centering
    \includegraphics[width=0.85\linewidth]{luminosidade.png}
    \caption{Variação da luminosidade entre 16:30h e 17:30h}
    \label{fig:luminosidade}
\end{figure}

Observa-se uma redução significativa da intensidade luminosa ao longo do período analisado, com valores reduzindo de aproximadamente 155 lux para cerca de 10 lux.

Esse comportamento está associado à diminuição gradual da iluminação natural no ambiente durante o final da tarde.

O gráfico evidencia o correto funcionamento do sensor BH1750 e da comunicação digital via protocolo I2C utilizada pelo sistema.

\subsection{Análise da Qualidade Relativa do Ar}

A Figura \ref{fig:qualidade_ar} apresenta os dados relativos à qualidade do ar obtidos pelo sensor MQ-135.

\begin{figure}[H]
    \centering
    \includegraphics[width=0.85\linewidth]{qualidade_do_ar.png}
    \caption{Variação da qualidade relativa do ar entre 16:30h e 17:30h}
    \label{fig:qualidade_ar}
\end{figure}

As leituras permaneceram em níveis baixos durante todo o período de aquisição, variando aproximadamente entre 37 e 61 unidades digitais do ADC.

Os valores registrados permaneceram dentro da faixa classificada pelo firmware como qualidade do ar ``BOA'', indicando baixa concentração relativa de gases detectáveis no ambiente monitorado.

Além disso, o comportamento do gráfico demonstra o funcionamento do processo de conversão analógico-digital realizado pelo ADC de 10 bits do microcontrolador. O sinal analógico fornecido pelo MQ-135 foi convertido em valores digitais entre 0 e 1023, permitindo que o firmware processasse as informações e classificasse a qualidade relativa do ar.

\subsection{Discussão Geral dos Resultados}

Os resultados obtidos demonstram o funcionamento integrado do sistema embarcado desenvolvido, evidenciando a comunicação entre sensores digitais e analógicos, processamento de dados e exibição das informações ambientais em tempo real.

O sistema apresentou estabilidade durante todo o período de aquisição, sem falhas de comunicação I2C ou inconsistências relevantes nas leituras realizadas.

Além disso, os resultados permitem observar na prática os principais elementos da arquitetura clássica de Von Neumann aplicados ao sistema desenvolvido:
\begin{itemize}
    \item os sensores atuam como dispositivos de entrada;
    \item o ATmega328P realiza o processamento dos dados;
    \item a memória interna armazena instruções e variáveis do firmware;
    \item o Monitor Serial atua como dispositivo de saída.
\end{itemize}

O projeto também demonstra a integração entre sinais digitais e analógicos dentro de um mesmo sistema computacional embarcado, especialmente através da utilização simultânea do protocolo I2C e do conversor analógico-digital interno do microcontrolador.

Dessa forma, os resultados experimentais validam o funcionamento da arquitetura implementada e demonstram a aplicação prática dos conceitos estudados em arquitetura e organização de computadores.

\section{Conclusão}

O presente trabalho permitiu o desenvolvimento e análise de um sistema embarcado de monitoramento ambiental utilizando sensores digitais e analógicos integrados ao microcontrolador ATmega328P.

Durante a execução do projeto, foi possível implementar a aquisição de dados ambientais em tempo real através dos sensores BME280, BH1750 e MQ-135, realizando medições de temperatura, umidade relativa do ar, pressão atmosférica, luminosidade e qualidade relativa do ar.

Os resultados experimentais obtidos demonstraram o funcionamento estável do sistema ao longo do período de monitoramento, validando tanto a comunicação I2C utilizada pelos sensores digitais quanto o processo de conversão analógico-digital aplicado ao MQ-135.

Além disso, o projeto possibilitou observar, na prática, diversos conceitos estudados em arquitetura e organização de computadores. A estrutura desenvolvida apresentou forte relação com a arquitetura clássica de Von Neumann, uma vez que o sistema embarcado implementa:
\begin{itemize}
    \item dispositivos de entrada, representados pelos sensores;
    \item processamento centralizado realizado pelo ATmega328P;
    \item armazenamento de instruções e variáveis em memória;
    \item dispositivos de saída, representados pelo Monitor Serial.
\end{itemize}

A integração entre sensores digitais e analógicos também demonstrou o funcionamento conjunto de diferentes formas de aquisição de dados dentro de um mesmo sistema computacional embarcado.

Outro aspecto relevante observado foi a atuação do ADC de 10 bits do ATmega328P, responsável pela conversão do sinal analógico contínuo do MQ-135 em valores digitais processáveis pelo firmware. Esse processo evidenciou a importância da conversão analógico-digital em sistemas computacionais modernos, permitindo que informações provenientes do meio físico sejam interpretadas e processadas digitalmente.

A utilização da linguagem Python para tratamento dos dados e geração dos gráficos também contribuiu para a análise temporal das variáveis ambientais, permitindo visualizar o comportamento do sistema ao longo do experimento realizado.

Dessa forma, conclui-se que o projeto atingiu os objetivos propostos, demonstrando de maneira prática a aplicação dos conceitos de arquitetura de computadores, sistemas embarcados, comunicação digital, conversão analógico-digital e integração hardware/software em um sistema funcional de monitoramento ambiental.

\renewcommand{\refname}{}
\section{Referências}
\begin{thebibliography}{9}

\bibitem{atmega}
ATMEL CORPORATION. \textit{ATmega328/P Datasheet Summary}. 
Atmel-42735A-ATmega328/P\_Datasheet\_Summary-06/2016, 2016.

\bibitem{bme280}
BOSCH SENSORTEC. \textit{BME280 Data Sheet: Digital Humidity, Pressure and Temperature Sensor}. 
Document revision 1.6, BST-BME280-DS002-15, September 2018.

\bibitem{bh1750}
ROHM SEMICONDUCTOR. \textit{BH1750FVI Ambient Light Sensor IC Series Datasheet}. 
ROHM Co., Ltd.

\bibitem{mq135}
HANWEI ELECTRONICS CO., LTD. \textit{MQ-135 Gas Sensor Technical Data}. 
MQ-135 Gas Sensor Datasheet.

\bibitem{tanenbaum}
TANENBAUM, A. S. \textit{Organização Estruturada de Computadores}. 6ª edição. São Paulo: Pearson, 2013.

\bibitem{stallings}
STALLINGS, William. \textit{Arquitetura e Organização de Computadores}. 8. ed. São Paulo: Pearson, 2010.

\end{thebibliography}
\end{document}
